\section{Introduction}
\label{sec:intro}
Selecting the right elements from a catalogue is a prerequisite for many language-model agents that access Web information. A tool-using agent must retain the operations needed to satisfy a request; a knowledge-grounded system must retain the entities that connect a question to its supporting facts. In either setting, choosing individually plausible names can omit intermediate context, while selecting too broadly can introduce irrelevant elements. Understanding selection reliability therefore requires examining how the requested elements compose.

Tool Graph Complexity reports that tool-selection success differs across dependency patterns, beyond what the number of reference tools alone describes~\cite{anonymousToolGraph}. Its controlled analysis retains a sequential-pattern penalty after accounting for measured size and surface properties. This motivates a question about the source of structural difficulty: does a pattern remain difficult when its relations are explicitly supplied, or does part of the difficulty arise from recovering dependencies from descriptions?

The same work investigates this question using a knowledge graph (KG) in Appendix~O. It reports partial structural transfer, an easier explicit-KG Chain profile, residual incoming-tree difficulty, and adjacent-decoy errors~\cite{anonymousToolGraph}. These antecedent observations motivate examining which errors remain after relation disclosure and how evidence conditions constrain their interpretation.

We address this question through an archived corporate-KG entity-selection study. The graph connects companies, people, products, locations, and news and was developed for corporate GraphRAG evaluation~\cite{kirpichev2026towards}. Our Web-related concern is the reliability of selecting connected evidence for querying interlinked corporate information: a route can be incomplete even when plausible endpoint entities are retained. We evaluate recovery of a declared entity-route reference, separating this access step from final answer correctness.

Our analysis distinguishes three comparisons. Raw versus Raw+E adds explicit relations to otherwise matched entity cards. Controlled triples versus templated prose and alternative card hosting retain designated facts while changing their presentation. Adjacent-context pairs add connected but query-irrelevant distractors. Keeping these contrasts separate matters: an information gain cannot be interpreted as merely an encoding effect, and complete route coverage is not the same objective as finding one sufficient answer.

We ask: \textbf{RQ1}, how do entity-selection profiles vary with task family and reference size? \textbf{RQ2}, how do relation availability, fact presentation, and adjacent contexts change recall and residual errors? \textbf{RQ3}, what do these observations imply for the interpretation of Tool-to-KG structural transfer? RQ3 is a descriptive comparison with prior evidence, not a new matched cross-domain experiment.

The additional contribution is a diagnostic follow-up: (i) separating relation disclosure from fact-preserving presentation, and (ii) localizing residual omissions, excess selection, and same-task presentation sensitivity. Family and size profiles provide its structural context. Recall is primary; precision and exact distinguish selectivity and complete-set recovery. This is an empirical analysis of an existing KG, without a new algorithm or resource claim. Exact task/response overlap with Appendix~O is unquantified; we do not claim independent replication.

The results show why relation visibility is insufficient as a complete explanation. Disclosure substantially increases recall, yet residual profiles differ across families and models. On six-node incoming trees, Qwen retains high precision but incomplete recall; repeating the same relations across both endpoint cards increases recall more than moving them to destination cards alone. Contrary model outcomes and regressions remain visible. These observations constrain possible explanations of selection errors without identifying a shared internal mechanism for tools and KG entities.
